
\section{One word against one: the tractable side and profile steps}\label{sec:tractable}

This section proves part~(i) of Theorem~\ref{thm:boundary}, the
fixed-profile criterion, the single-word claims of
Theorem~\ref{thm:directed}, and the P-complete half of
Corollary~\ref{cor:algebraic}.
% gloss: single-word class
We call a model \defterm{single-word}, or in the \defterm{single-word class}, when all its word tests are word-against-word.

% ---------------------------------------------------------------------------
\subsection{From word tests to binomial equations}\label{sec:compile}

Feasibility then splits into three systems with separate variables.

\begin{lemma}[Feasibility separates]\label{lem:separate}
Fix a shared-probe model in the single-word class, a profile $p$ whose held
word tests form an admissible set, and a parameter vector
$(\theta,\alpha,L)$ with tables~$P$. Then the configuration $(p,P)$ is
feasible if and only if conditions (a)--(c) hold.
(a)~The vector $\theta$ satisfies $\theta(w)=\theta(w')$ for every held
word test $w\sim w'$ and $\theta(w)=c\,\theta(w')$ for every held
probability-ratio test.
(b)~The vector $\alpha$ is a feasible point of $\Phi_p$ with its budgets.
(c)~The tables satisfy $L_m(g,h;r,r')>0$ exactly when
$(m,g,r,h,r')\in\Gamma_p$, and $L$ satisfies the block-mass conditions at
every level of~$p$.
\end{lemma}

Lemma~\ref{lem:compile}(a) gives the split
(Appendix~\ref{app:tractable-proofs}). The next lemma moves between steps in the tables and
steps in each part.

\begin{lemma}[Tables and parameters]\label{lem:transfer}
Fix a shared-probe model and put $N=|X||R|$.
\begin{enumerate}
\item[(a)] Any two parameter vectors, with eigenvalue vectors
$\alpha_A,\alpha_B$ and tables $P_A,P_B$, satisfy
$\|\alpha_A-\alpha_B\|_\infty\le(2N/\delta)\,\|P_A-P_B\|_\infty$.
\item[(b)] On the set of parameter vectors, the tables are uniformly
continuous functions of $(\theta,\alpha,L)$.
\item[(c)] At a profile $p$ of the single-word class, let rational vectors
$\theta,\theta'$ satisfy the probability floor, normalization and the
equations of Lemma~\ref{lem:separate}(a). For every rational
$\varepsilon>0$, a finite sequence of such rational vectors with sup-norm
steps below~$\varepsilon$ joins them.
\item[(d)] The same holds for two register tables satisfying
Definition~\ref{def:model} and Lemma~\ref{lem:separate}(c) at~$p$, through
rational affine interpolation.
\end{enumerate}
\end{lemma}

Appendix~\ref{app:tractable-proofs} proves the lemma, interpolating label
probabilities geometrically and registers affinely.

\begin{remark}[The role of each restriction]\label{rem:restrictions}
The $\ell^1$ budget keeps every base kernel positive along the monomial
curves of Lemma~\ref{lem:restriction}(a), which shrink absolute values,
while the eigenvalues themselves can be negative or zero. Positivity of
$Q_m$ alone would not do. On $(\mathbb Z/2)^2$, take two generating
characters and their product as the three probes. Eigenvalues
$(4/5,4/5,7/10)$ give a positive kernel. Damping the third eigenvalue to
zero gives the entry $(1-4/5-4/5)/4=-3/20$. The probability floor
$\theta_m\ge\delta>0$ gives Lemma~\ref{lem:transfer}(a), through division
by~$\theta_m$. Equal word lengths make the rate equations homogeneous, so
normalization preserves them. The register conditions are affine and
separate from the others.
\end{remark}

% ---------------------------------------------------------------------------
\subsection{Proof of the fixed-profile criterion}\label{sec:fixed-proof}

\begin{proof}[Proof of Theorem~\ref{thm:fixed-profile}]
By Lemma~\ref{lem:separate}, the eigenvalue vectors of feasible
configurations at~$p$ are exactly the feasible rational points of $\Phi_p$
with its budgets. The reversal equations belong to $\Phi_p$ because every
parameter vector satisfies them. Every coordinate lies in a budget, so this
feasible rational set is bounded, and Theorem~\ref{thm:core} applies to it.

For necessity, fix a rational $\varepsilon_\alpha>0$ and a reconfiguration
sequence from $A$ to $B$ at tolerance $\delta\varepsilon_\alpha/(2N+1)$,
with $N$ as in Lemma~\ref{lem:transfer}. By Lemma~\ref{lem:transfer}(a),
consecutive eigenvalue vectors differ by less than $\varepsilon_\alpha$ in
sup norm. So $\alpha_A$ and $\alpha_B$ are finely connected, and
Theorem~\ref{thm:core} gives $\SC_p(\alpha_A)=\SC_p(\alpha_B)$.

For sufficiency, fix $\varepsilon>0$ and change one part of the parameter
vector at a time. Lemma~\ref{lem:transfer}(c) joins $\theta_A$ to
$\theta_B$, Theorem~\ref{thm:core} joins $\alpha_A$ to $\alpha_B$, and
Lemma~\ref{lem:transfer}(d) joins $L_A$ to $L_B$, whose positive entries
both give~$\Gamma_p$. Every intermediate configuration is feasible by
Lemma~\ref{lem:separate}. By Lemma~\ref{lem:transfer}(b), fine enough steps
in each part change every table entry by less than~$\varepsilon$.
\end{proof}

% ---------------------------------------------------------------------------
\subsection{The algorithm and its certificates}\label{sec:algorithm}

The linear-programming (LP) computation of Section~\ref{sec:computing-core}
has at most $|J|$ deletion rounds and one final check. Its LPs have
polynomial input size because their matrices are restrictions of the
binary-encoded exponent differences. Neither this computation nor the real
signed-core criterion of Corollary~\ref{cor:real} requires a bounded locus,
so the same computation decides real path-component membership
(Corollary~\ref{cor:real-complexity}). Fraca and Haddad's
Proposition~27~\cite{fracahaddad} analyzes their Algorithm~2
(Section~\ref{sec:computing-core}).

% gloss: Gordan certificate
A peeling sequence shows which coordinates can be removed; a certificate
must also show that no more can be. A \defterm{Gordan certificate} for a
support $C$ is an integer vector $\lambda$, indexed by the tests active
on~$C$, with $\Delta_C^{\mathsf T}\lambda>0$ in every coordinate; for
$C=\emptyset$ it is empty. In the sign example on $C=\{y,z\}$, multiplier $1$ on
the difference vector $(2,-2)$ of $y^2=z^2$, multiplier $3$ on $(0,1)$ of
$z=1/4$, and $0$ on the identity reversal equations give
$\Delta_C^{\mathsf T}\lambda=(2,1)$. Let $D_0$ be the
maximum of $1$ and the absolute values of the entries of the vectors $a-b$.

\begin{theorem}[Polynomial time, with certificates]\label{thm:certificates}
Let $\Phi$ be an explicitly binary-encoded budgeted signed-binomial system
with $n=|J|$ coordinates and $|\Phi|$
tests.
\begin{enumerate}
\item[(a)] For every support $S$, the core $\core_\Phi(S)$ is computable in
deterministic polynomial time.
\item[(b)] Integer weight vectors of a peeling sequence from $S$ to~$C$,
each checked against the tests active at its own support, together with a
Gordan certificate for $C$, prove $C=\core_\Phi(S)$. Such a certificate
exists with at most $n$ stages and at most $n$ nonzero multipliers, each
of absolute value at most $n!\,D_0^n$. It can be stored with binary length
$O(n^3(1+\log(n+1)+\log D_0)+n\log(|\Phi|+1))$.
% gloss: polynomial-size certificates (G8; read back after the proof)
\item[(c)] Fine Reachability at a fixed profile in the single-word class is
decided in deterministic polynomial time, and both answers have
polynomial-size certificates, checkable in polynomial time.
\end{enumerate}
\end{theorem}

\begin{proof}
Part~(a) is the computation above. Appendix~\ref{app:tractable-proofs}
proves the bounds and~(b)--(c). In~(b), soundness uses the terminal
dot-product check and completeness uses Gordan's alternative. Part~(c)
compiles $\Phi_p$, applies Theorem~\ref{thm:fixed-profile} and gives the
verifier for both verdicts.
\end{proof}

Thus both answers are certified by peeling stages and Gordan certificates,
not by the points of a sequence.

Checking whether a supplied rational point satisfies the binomial equations
of a Fine Connectivity instance is a separate task. After handling zero
coefficients, vanishing monomials and signs, clear denominators in each
comparison, keeping the powers encoded. Proposition~2.1 of Etessami,
Stewart and
Yannakakis~\cite[Proposition~2.1]{etessamistewartyannakakis} supplies the
power-product equality check. With the zero and sign checks and
rational budget arithmetic, this equality check verifies feasibility of the
supplied endpoint. Boundedness remains a promise, and the check does not
decide whether a feasible point exists.

\begin{proposition}[Explicit seq{}uences can need exponential binary length]\label{prop:long-sequences}
For $k\ge1$, let $\Phi$ have coordinates $x_0,\dots,x_k$, tests
$x_{i+1}=x_i^2$ for $0\le i<k$, and budgets $|x_i|\le1$. The points $0$
and $(1,\dots,1)$ are finely connected. Every sequence of feasible rational
points between them with sup-norm steps below $1/4$ contains a point whose
coordinate $x_k$ has a reduced denominator of binary length at least
$2^k+1$.
\end{proposition}

The proof (Appendix~\ref{app:tractable-proofs}) follows~$x_0$ across
$[1/4,1/2)$.

\paragraph{Scope of the tolerance.} The single-word P-completeness result
and its certificates concern fine reachability, at every positive rational
tolerance. Different signed cores can still be joined at a coarse
tolerance, as one step joins the two signs of the sign example above~$1/8$. Proposition~\ref{prop:separation} supplies a
sufficient separation cutoff whose expanded rational denominator can
require exponentially many binary digits. For each NCL reduction instance
of Section~\ref{sec:hard}, the NCL verdict agrees with sequence existence
at its supplied cutoff and at every smaller positive tolerance. A positive
certificate proves sequence existence at every tolerance, while by
Proposition~\ref{prop:long-sequences} an explicit sequence can require
exponentially many binary digits even at a fixed tolerance.

% ---------------------------------------------------------------------------
\subsection{P-hardness}\label{sec:p-hardness}

% gloss: anchor
As in both examples of Section~\ref{sec:question}, $\mathsf z$ is the
\defterm{anchor}, a label calibrated to $1/4$ whose coordinate never peels.
We reduce from monotone circuit value, which Goldschlager proved P-complete
under log-space reductions~\cite{goldschlager}; Greenlaw, Hoover and Ruzzo
prove the variant with numbered gates and assigned
inputs~\cite[Theorem~6.2.2]{greenlawhooverruzzo}. A circuit lists gates
$1,\dots,G$, each a constant $0$ or $1$ or a binary AND or OR of earlier
gates, and names an output gate~$\gamma_{\mathrm{out}}$.

% gloss: sentinel
The reduction outputs an instance of the single-word class with base
$\{0,1\}$, one probe, a one-value register and only level~$0$. Its labels,
each its own reversal, are a gate label $\mathsf g_\gamma$ for each gate,
the anchor and a \defterm{sentinel}~$\mathsf s$, whose sign alone
distinguishes the two endpoints. The eigenvalues of the gate labels, the anchor and the sentinel are
$g_\gamma$, $z$ and~$s$, respectively. With $\delta=1/(G+2)$ every label
probability is~$\delta$; also $\kappa=1/2$ and $H=3$. To keep two tests per
gate uniformly for log-space output, each gate receives two named word tests,
its \emph{slots}. An AND gate uses one test per input, constant~$0$ and OR
repeat one test, and constant~$1$ uses two trivial identities so it still
fills both slots. With inputs $\gamma',\gamma''$, the slots are
\[
\begin{array}{lll}
\text{constant }0: & \mathsf g_\gamma\mathsf{zz}\sim\mathsf{zzz}, & \mathsf g_\gamma\mathsf{zz}\sim\mathsf{zzz};\\
\text{constant }1: & \mathsf{zzz}\sim\mathsf{zzz}, & \mathsf{zzz}\sim\mathsf{zzz};\\
\text{AND}: & \mathsf g_\gamma\mathsf g_{\gamma'}\mathsf z\sim\mathsf z\mathsf g_{\gamma'}\mathsf z, &
\mathsf g_\gamma\mathsf g_{\gamma''}\mathsf z\sim\mathsf z\mathsf g_{\gamma''}\mathsf z;\\
\text{OR}: & \mathsf g_\gamma\mathsf g_{\gamma'}\mathsf g_{\gamma''}\sim\mathsf z\mathsf g_{\gamma'}\mathsf g_{\gamma''}, &
\mathsf g_\gamma\mathsf g_{\gamma'}\mathsf g_{\gamma''}\sim\mathsf z\mathsf g_{\gamma'}\mathsf g_{\gamma''}.
\end{array}
\]
The output test is
$\mathsf{ss}\mathsf g_{\gamma_{\mathrm{out}}}\sim\mathsf{zz}\mathsf g_{\gamma_{\mathrm{out}}}$.
All $2G+1$ word tests and the calibration $z=1/4$ are held, and the only
admissible set is the set of all word tests. The endpoint $A$ has every
eigenvalue $1/4$, and $B$ differs only in $s=-1/4$; both are feasible
whatever the circuit, since $s$ enters only the output test, squared.

In other words, a test is active only while all its letters are nonzero.
So the coordinates that peel follow gate-local Horn rules: a constant~$1$
is removed, an AND gate once both inputs are, an OR gate once either is,
and $s$ once the output gate is. In the three-gate circuit, $\mathsf g_1$
is the constant~$0$, $\mathsf g_2$ the constant~$1$, and the output is
$\mathsf g_3=\mathrm{OR}(\mathsf g_1,\mathsf g_2)$. Here $g_2$ peels first,
making the test of $\mathsf g_3$ inactive; then $g_3$ peels, making the
output test inactive; then $s$ peels. The core is $\{z,g_1\}$, and
$B$ is finely reachable from~$A$. With AND in place of OR, the slot of
$\mathsf g_1$ keeps $g_3=z$ after $g_2$ leaves. Then $\{z,g_1,g_3,s\}$ is
irreducible and $B$ is not finely reachable, although one numerical step
joins the endpoints at any tolerance above $1/20$.

\begin{theorem}[P-hardness]\label{thm:p-hard}
The map from a monotone circuit to the instance above is computable in
deterministic logarithmic space. The configurations $A$ and $B$ are
feasible at one common profile with full table support, and $B$ is finely
reachable from $A$ if and only if the circuit outputs~$1$. Hence Fine
Reachability at a fixed profile in the single-word class is P-hard under
deterministic log-space many-one reductions on promised-feasible inputs.
\end{theorem}

\begin{proof}[Proof sketch]
An active gate test other
than a constant-$1$ slot gives $\eta_{g_\gamma}=\eta_z$, where $\gamma$ is
its gate. The constant-$1$ slots are identities. The active output test
gives $2\eta_s=2\eta_z$, and the calibration gives $\eta_z=0$. The anchor,
the false gates and, if the output is~$0$, the sentinel form an irreducible
support, so they lie in the core of the full support
(Lemma~\ref{lem:irreducible}). A first true gate in that core would peel
alone, and so would the sentinel once the output gate is gone. So the
signed cores of $A$ and~$B$ agree exactly when the output is~$1$, and
Theorem~\ref{thm:fixed-profile} gives the equivalence.
Appendix~\ref{app:tractable-proofs} gives the details and the log-space
transducer.
\end{proof}

Thus the sentinel survives peeling exactly when the circuit outputs~$0$.
With
Theorem~\ref{thm:certificates}, the problem is P-complete, and since the
reductions land inside the promise, it is in log-space-uniform NC only if P
equals log-space-uniform NC.

The reduction follows the circuit-value pattern of Fraca and Haddad's
Proposition~32 for continuous Petri nets~\cite{fracahaddad}. Here the
surviving signed core computes the circuit. For a given gate count, the
endpoint eigenvalue vectors are fixed independently of the gate types,
wiring and output value; the endpoints share a profile within each
instance and differ only in one sign. On these instances the removed set
is the least model of the gate-local Horn rules above, which the
propagation of Dowling and Gallier computes in linear
time~\cite[Sections~4--5]{dowlinggallier}. This reading is for the circuit
instances only: in general a peeling step combines tests through kernel
vectors found by linear programming.

% ---------------------------------------------------------------------------
\subsection{Profile steps: the directed criterion}\label{sec:directed}

% gloss: profile order
% gloss: cover
% gloss: profile step
Theorem~\ref{thm:directed} lets the profile change: a configuration can
acquire held tests, refine its register partitions or shrink its table
support, but never the reverse, so a held test keeps passing for the rest
of the sequence. Definition~\ref{def:profile-steps} extends the sequences and fine
reachability of Section~\ref{sec:profiles}.

\begin{definition}[Profile order and profile steps]\label{def:profile-steps}
For profiles $p$ and~$q$, put $p\preceq q$ if $\Gamma_q\subseteq\Gamma_p$,
$F^q_\ell$ refines $F^p_\ell$ at every level~$\ell$, and
$\mathcal T_p\subseteq\mathcal T_q$. An attained profile $q$ \emph{covers}
an attained $p$ if $p\prec q$ and no attained profile lies strictly between.
A \emph{numerical step} at tolerance~$\varepsilon$ joins feasible
configurations $A,B$ with $p_A=p_B$ and $\|P_A-P_B\|_\infty<\varepsilon$.
A \emph{profile step} joins feasible configurations $A,B$ such that $p_B$
covers $p_A$ and either $\Gamma_{p_B}=\Gamma_{p_A}$ and $P_B=P_A$, or
$\Gamma_{p_B}\subsetneq\Gamma_{p_A}$ and the projected tables
$\theta_mQ_m(g,h)$ agree. A \emph{reconfiguration sequence at
tolerance~$\varepsilon$} is a finite, possibly empty, sequence of such
steps through feasible rational configurations. $B$ is \emph{finely
reachable} from $A$ if one exists from $A$ to $B$ for every rational
$\varepsilon>0$.
\end{definition}

In other words, steps go up the partial order~$\preceq$, and agreement of
the projected tables means $\theta_B=\theta_A$ and $\alpha_B=\alpha_A$, so
only the register can change. %
In the sign example the two labels are self-reversing. Take $\kappa=1/2$,
$H=2$, full table support and only level~$0$. The admissible sets of word tests are the empty set and the
singleton containing $\mathsf{yy}\sim\mathsf{zz}$. Let $A$ hold only
the calibration, at $y=1/4$, and let $B$ also hold
$\mathsf{yy}\sim\mathsf{zz}$, at $y=-1/4$; then $p_A\preceq p_B$.

\begin{center}
\fbox{\begin{minipage}{0.93\linewidth}
\textsc{Fine Reachability with profile steps}\par\smallskip
\emph{Input:} an explicitly presented model, and two rational configurations
$A$ and~$B$, each given by its profile, or data that determine it
(Remark~\ref{rem:readback}), and its parameter vector or transition tables,
in binary.\par\smallskip
\emph{Promise:} $A$ and $B$ are feasible.\par\smallskip
\emph{Question:} is $B$ finely reachable from~$A$?\par\smallskip
\emph{Supplied-tolerance variant:} the input also contains
$\varepsilon_0\in\QQ_{>0}$, and the question is whether a reconfiguration
sequence at tolerance~$\varepsilon_0$ exists.
\end{minipage}}
\end{center}

\begin{lemma}[One change per profile step]\label{lem:covers}
If $q$ covers $p$, then $q$ differs from $p$ in exactly one of the
following: the table support; the held word tests; the held value tests;
the partition at one level. If the held value tests differ, $q$ holds
exactly one more value test than~$p$.
\end{lemma}

Appendix~\ref{app:tractable-proofs} proves the lemma by building an attained
profile, above $p$ and at most~$q$, that changes one part. So every step has one of six types: numerical; support;
word-test; value-test; partition at a non-final level; partition at the
final level. Only a word-test cover can add more than one test. In that
case no admissible set of word tests lies strictly between the two held
word-test sets.

\begin{theorem}[Directed criterion]\label{thm:directed}
Fix a shared-probe model in the single-word class, and let $A$ and $B$ be
feasible rational configurations. Then $B$ is finely reachable from $A$ if
and only if $p_A\preceq p_B$ and $\SC_{p_A}(\alpha_A)=\SC_{p_A}(\alpha_B)$,
where both signed cores are computed from the eigenvalue constraints
$\Phi_{p_A}$ of the source profile, including the reversal equations. Fine
Reachability with profile steps in this class is P-complete under
deterministic log-space many-one reductions, and both answers have
polynomial-size certificates, checkable in polynomial time, that include
the order comparison. For models in which exactly one word test is
word-against-pair and all others are word-against-word, Fine Reachability
with profile steps is PSPACE-hard under polynomial-time many-one
reductions.
\end{theorem}

Thus profile steps add only the order condition $p_A\preceq p_B$, and both
cores come from the source's constraints; the sign example below shows why
the target's constraints would not do.

\begin{proof}[Proof outline]
Steps never lower the profile, so every later configuration passes the
tests of~$p_A$, and profile steps preserve~$\alpha$; necessity then follows as
for Theorem~\ref{thm:fixed-profile}. For sufficiency, combine the target's
label probabilities and eigenvalues with the source register at~$p_A$, and
reach that configuration by Theorem~\ref{thm:fixed-profile}. Then climb a
saturated chain of attained profiles at these values, interpolating to the
upper register at the lower profile before each cover with unchanged
support. Appendix~\ref{app:tractable-proofs} gives the details and proves
the complexity claims.
\end{proof}

In the sign example, $\Phi_{p_A}$ gives only $\eta_z=0$, so $y$ peels and
both signed cores are $z$ alone: $B$ is finely reachable from~$A$, through
$y=0$. At this calibration-only profile, every feasible eigenvalue vector
has signed core $\{z\}$ with positive sign. Thus every feasible target at
or above that profile is finely reachable from any feasible source there,
by Theorem~\ref{thm:directed}. Using the target's constraints for both cores would add
$2\eta_y=2\eta_z=0$, incorrectly ruling out fine reachability from $A$ to
$B$. Let $A^+$ hold the test at the parameter vector of~$A$. From $A^+$ the
cores differ in the sign of~$y$, so $B$ is not finely reachable; the exact
threshold $1/8$ of Section~\ref{sec:question} is for this pair. No
sequence leads from $A^+$ back to $A$ at any tolerance: reachability is
between configurations, not parameter vectors.

The same sign obstruction persists in a model with a nontrivial register and
four labels. On the staged example's base, register values and labels, take
label probabilities $1/4$ and the common full-support profile holding
$\mathsf{yy}\sim\mathsf{zz}$ and $z=1/4$. Take eigenvalues
$(x,y,z,f)=(-1/8,\pm1/4,1/4,1/8)$, the register redrawn from $\nu$ at every
step, and one block at every level. Then Lemma~\ref{lem:transfer}(a), with
$N=8$ and $\delta=1/8$, excludes every sequence between the two signs at
tolerance $1/272$, a sufficient tolerance rather than the exact threshold.
At the same profile, a direct step between these endpoints changes a table
entry by at most $1/64$, so it is allowed at any tolerance above $1/64$.

Both conditions of Theorem~\ref{thm:directed} can be evaluated from the
full-state labeled transition tables, decision availability at each level,
and held-test membership recorded in Remark~\ref{rem:readback}
(Appendix~\ref{app:tractable-proofs}).

The ordinary and signed row sums in Lemma~\ref{lem:transfer}(a) give $\theta_m$ and $\theta_m\alpha_{mi}$. Dividing by the positive label probability yields $\alpha_{mi}$, so the tables determine both types of value test; Appendix~\ref{app:tractable-proofs}
reconstructs the rest of the configuration.

% ---------------------------------------------------------------------------
\subsection{The algebraic form, first half, and the circuit transfer}\label{sec:circuit-transfer}

Section~\ref{sec:algorithm} and Theorem~\ref{thm:core} give the upper
bound in Corollary~\ref{cor:algebraic}, under the promise that $\FQ$ is
bounded. For hardness, the transducer of Section~\ref{sec:p-hardness} emits
only the eigenvalue tests, budgets $|x_m|\le1/2$ and endpoint
eigenvalue vectors; the resulting instances meet the promises
(Appendix~\ref{app:tractable-proofs}). By Corollary~\ref{cor:real}, which
needs no boundedness, these instances also give the P-hardness in
Corollary~\ref{cor:real-complexity}.
